\worksheettitle{Условия между цифрами}{\modulemark{M07} Учебный модуль \textbullet\ упорядоченный перебор}
\studentnameblock

\begin{checkpointbox}[Вход после M05]
  Из цифр 1, 3, 5 составьте все двузначные числа без повторений:
  \answerblank[85mm]. Как вы убедились, что список полный?
  \answerblank[80mm].
\end{checkpointbox}

\section{Переводим условие на язык цифр}
Пусть (a) — цифра десятков, (b) — цифра единиц. Тогда двузначное число
имеет вид (10a+b), причём (a) принимает значения от 1 до 9, а (b) — от
0 до 9.

\begin{fillbox}[Единицы на 2 больше десятков]
  Фраза означает: к цифре десятков прибавили 2 и получили цифру единиц.
  Короткая запись:
  [
    b=a+2.
  ]
  Если $a=1$, то $b=$\answerblank[12mm], число \answerblank[20mm].
  Если $a=2$, то $b=$\answerblank[12mm], число \answerblank[20mm].
\end{fillbox}

\begin{warningbox}[Не обращайте зависимость]
  «Единицы на 2 больше десятков» — это (b=a+2), а не (a=b+2).
  Сначала назовите, какая цифра больше, и только затем записывайте равенство.
\end{warningbox}

\section{Строим полный список}
Заполняйте таблицу, увеличивая (a) на 1.

\begin{center}
\renewcommand{\arraystretch}{1.45}
\begin{tabular}{c|ccccccccc}
  \toprule
  (a) & 1 & 2 & 3 & 4 & 5 & 6 & 7 & 8 & 9 \\
  \midrule
  (b=a+2) & 3 & 4 & \answerblank[8mm] & \answerblank[8mm] &
  \answerblank[8mm] & \answerblank[8mm] & \answerblank[8mm] & 10 & 11 \\
  \bottomrule
\end{tabular}
\end{center}

\begin{reflectionbox}[Граница перебора]
  Значения 10 и 11 не являются цифрами. Поэтому последняя допустимая пара —
  $a=$\answerblank[12mm], $b=$\answerblank[12mm], а полный список чисел:
  \answerblank[105mm].
\end{reflectionbox}

\newpage
\section{Другие виды условий}

\begin{fillbox}[Единицы в 2 раза больше десятков]
  Условие: (b=2a). Заполните допустимые пары.
  \begin{center}
  \renewcommand{\arraystretch}{1.4}
  \begin{tabular}{c|ccccc}
    \toprule
    (a) & 1 & 2 & 3 & 4 & 5 \\
    \midrule
    (b=2a) & 2 & \answerblank[10mm] & \answerblank[10mm] &
      \answerblank[10mm] & 10 \\
    \bottomrule
  \end{tabular}
  \end{center}
  Подходящие числа: \answerblank[80mm]. Почему перебор заканчивается?
  \answerblank[75mm].
\end{fillbox}

\begin{fillbox}[Сумма цифр равна 7]
  Условие: (a+b=7), значит (b=7-a).
  \begin{center}
  \renewcommand{\arraystretch}{1.4}
  \begin{tabular}{c|ccccccc}
    \toprule
    (a) & 1 & 2 & 3 & 4 & 5 & 6 & 7 \\
    \midrule
    (b) & 6 & \answerblank[8mm] & \answerblank[8mm] & \answerblank[8mm] &
    \answerblank[8mm] & \answerblank[8mm] & 0 \\
    \bottomrule
  \end{tabular}
  \end{center}
  Полный список: \answerblank[100mm].
\end{fillbox}

\begin{practicebox}[\levelbase. Сначала запишите равенство]
  \begin{enumerate}[label=\textbf{\arabic*.}]
    \item Единицы на 4 больше десятков: \answerblank[35mm].
    \item Десятки на 3 больше единиц: \answerblank[35mm].
    \item Единицы в 3 раза больше десятков: \answerblank[35mm].
    \item Сумма цифр равна 9: \answerblank[35mm].
  \end{enumerate}
\end{practicebox}

\newpage
\section{Практика: выбирает учитель}

Для каждого задания запишите зависимость, постройте таблицу или группы и
объясните границу перебора.

\begin{practicebox}[\levelpractice. Найдите все двузначные числа]
  \begin{enumerate}[label=\textbf{\arabic*.}]
    \item Цифра единиц на 1 больше цифры десятков.
    \item Цифра единиц на 3 меньше цифры десятков.
    \item Цифра единиц в 3 раза больше цифры десятков.
    \item Цифра десятков в 2 раза больше цифры единиц.
    \item Сумма цифр равна 5.
    \item Сумма цифр равна 10.
  \end{enumerate}
  \vspace{42mm}
\end{practicebox}

\begin{practicebox}[Проверьте количество]
  Для каждого предыдущего списка укажите число вариантов:
  \begin{center}
    1) \answerblank[12mm]\quad
    2) \answerblank[12mm]\quad
    3) \answerblank[12mm]\quad
    4) \answerblank[12mm]\quad
    5) \answerblank[12mm]\quad
    6) \answerblank[12mm].
  \end{center}
\end{practicebox}

\newpage
\section{Разбираемся в ошибках}

\begin{warningbox}[Исправьте решения]
  \begin{enumerate}[label=\textbf{\arabic*.}]
    \item «Единицы на 2 меньше десятков» записали (b=a+2).
    \item При условии (b=a+4) включили число 6,10.
    \item Для суммы цифр 6 записали (15,24,33,42,51), но забыли 60.
    \item Для условия (a=2b) записали (12,24,36,48).
  \end{enumerate}
  Для каждого пункта назовите причину ошибки и дайте правильный ответ.
  \vspace{38mm}
\end{warningbox}

\begin{fillbox}[Почему нельзя остановиться раньше]
  Ученик нашёл числа (14,25,36) для условия (b=a+3) и решил, что этого
  достаточно. Какие варианты он потерял? \answerblank[70mm].
  Какая таблица или закономерность доказывает полноту?
  \answerblank[100mm].
\end{fillbox}

\newpage
\section{Смешанные условия}

\begin{practicebox}[Сначала отберите, затем проверьте]
  \begin{enumerate}[label=\textbf{\arabic*.}]
    \item Найдите двузначные числа, у которых сумма цифр 8, а цифра десятков
      больше цифры единиц.
    \item Найдите двузначные числа, у которых единицы на 2 больше десятков и
      число больше 40.
    \item Найдите двузначные числа, у которых десятки в 2 раза больше единиц,
      а число меньше 80.
  \end{enumerate}
  \vspace{30mm}
\end{practicebox}

\begin{fillbox}[Составьте собственное условие]
  Придумайте условие между цифрами, которому удовлетворяют хотя бы три, но не
  более шести двузначных чисел.

  Условие словами: \answerblank[115mm].

  Равенство: \answerblank[35mm]. Список: \answerblank[55mm].

  Почему список полный: \answerblank[80mm].
\end{fillbox}

\newpage
\section{Контрольная точка}
\begin{checkpointbox}[Выполните без подсказки]
  Найдите все двузначные числа, в которых цифра единиц на 3 больше цифры
  десятков.
  \begin{enumerate}[label=\textbf{\arabic*.}]
    \item Запишите условие равенством: \answerblank[45mm].
    \item Полный список: \answerblank[70mm].
    \item Почему перебор заканчивается именно здесь?
      \answerblank[75mm].
  \end{enumerate}
\end{checkpointbox}

\begin{reflectionbox}[Моя готовность]
  \choicebox\ Не обращаю зависимость между цифрами.

  \choicebox\ Помню, что цифра принимает значения только от 0 до 9.

  \choicebox\ Изменяю одну цифру по порядку и вычисляю вторую.

  \choicebox\ Объясняю, почему список полный.

  \choicebox\ Мне нужна таблица M07-R.
\end{reflectionbox}

\begin{notebox}[Если закончили раньше]
  Выполните M07\(\star\): подсчитайте число вариантов без записи полного
  списка и свяжите формулу с допустимыми строками таблицы.
\end{notebox}
